\section{Strict ordinary flags over integral test algebras}
\label{sec:cubic-integral-ordinary}

The identities of AC-G57 admit an integral strengthening.  We
retain the actual good-prime hypotheses $\ell\ne2,3,269$ and
\texttt{GD.ine}, \texttt{GD.gen}: on inertia $\psi^2=1$ and
$\ell-1\nmid2e$.  Put $\gamma=\psi\kappa_\ell^e$ and
$(\gamma_0,\gamma_1,\gamma_2)=(\gamma,1,\gamma^{-1})$.
The cubic character is unramified at these primes.  The
multiplier $\eta$ is unramified on the local group.
All assertions below concern the actual local models of
AC-G49 and AC-G54; their construction and published local
inputs remain hypotheses exactly as stated there.

\begin{theorem}[AC-G57-integral: ordinary inertia over every local test algebra]
\label{AC-G57-integral}
Let $A$ be a commutative local $O$-algebra whose residue map
agrees with the fixed residual representation, and let $r$
be the specialization to $A$ of either actual local model.
There is a matrix $h\in\mathrm{GL}_3(A)$ with
$h\equiv I\pmod{\mathfrak m_A}$ such that the matrices of
$h^{-1}rh$ are upper triangular on $\ker\eta$, and on
inertia have ordered diagonal
$\gamma,1,\gamma^{-1}$.
This is the exact strict integral condition
\texttt{UnitaryThreeInertia.CondT}, defined by
\texttt{InHat} and \texttt{TriIn} in
\path{Development/AutomorphicArithmetic/AutomorphicVarieties/UnitaryThreeDeformations/AdjointDiagonalRepPreservesWeightPiece.lean}; it applies in particular
to the characteristic-$\ell$ complete DVR specializations
used by the source finiteness criterion.  Neither
characteristic zero, reducedness, nor passage to a
fraction field is required for this assertion.
\end{theorem}
\begin{proof}
Conjugation here is conjugation in the semidirect group;
on $\ker\eta$ it is ordinary matrix conjugation.
First we justify the polynomial identities in both models.
For the new split model they are (G57.4).  For the transported
model, use
\path{Development/AutomorphicArithmetic/AutomorphicVarieties/PadicHodgeTheory/ScalarMap.lean},
\texttt{UpperTheory.V.Loc.lemma2a}.  Its displayed statement
quantifies over an arbitrary commutative $O$-algebra $C$;
only $O$, not $C$, is assumed to have characteristic zero.
The premise \texttt{D.On U r} means that the representation
is obtained by a map from $U.R$ killing the actual ideal
$I_t$.  Thus it applies directly to the canonical map
$U.R\to U.R/I_t$.  Its conclusion is the characteristic
polynomial identity and the ordered three-factor product
identity in (G57.4).  This source theorem retains its
\texttt{LocalPublishedTheorems} premise, already retained
in AC-G49.  Transport by the exact conjugation of AC-G49
preserves both identities: the multiplier is one on
inertia, so there the semidirect conjugation is ordinary
matrix conjugation.  In both cases these polynomial
identities survive every $O$-algebra specialization to $A$.

Choose $s\in I_\ell$ with
$\overline\kappa_\ell(s)^{2e}\ne1$.  This is possible
because the mod-$\ell$ cyclotomic character on rational
inertia is onto $\mathbb F_\ell^\times$ and
$\ell-1\nmid2e$.  Since $\overline\psi(s)^2=1$,
the values
\[
 d_0=\gamma(s),\qquad d_1=1,\qquad d_2=\gamma(s)^{-1}
 \tag{G57I.1}
\]
are pairwise distinct modulo $\mathfrak m_A$.
Consequently every $d_i-d_j$, $i\ne j$, is a unit in $A$.
Write $M=r(s)$.  Its characteristic polynomial is
$\prod_i(X-d_i)$, so Cayley--Hamilton and Lagrange
interpolation give orthogonal idempotents
\[
 E_i=\prod_{k\ne i}\frac{M-d_k I}{d_i-d_k},\qquad
 \sum_iE_i=I,\qquad ME_i=d_iE_i.                 \tag{G57I.2}
\]
These equalities hold over $A$: divide polynomial identities
by the monic polynomial $\prod_i(X-d_i)$, or apply the
Chinese remainder theorem to the pairwise comaximal ideals
$(X-d_i)$ in $A[X]$.  Reducing modulo $\mathfrak m_A$, the
fixed residual inertia representation is diagonal with
these distinct entries.  Hence $\overline E_i$ is the
$i$th coordinate projector.  If $e_0,e_1,e_2$ is the
standard basis, the matrix
\[
             h=(E_0e_0,\ E_1e_1,\ E_2e_2)       \tag{G57I.3}
\]
reduces to $I$.  Its determinant is a unit because $A$ is
local, so $h$ is invertible and $h^{-1}Mh=D$ with
$D=\operatorname{diag}(d_0,d_1,d_2)$.

Work in this basis.  For $x\in I_\ell$, set $N=r(x)$
and $p_i=\gamma_i(x)$.  The ordered product identity,
at $(x,s,s)$, $(s,x,s)$ and $(s,s,x)$, gives respectively
\[
\begin{split}
 (N-p_0I)(D-d_1I)(D-d_2I)&=0,\\
 (D-d_0I)(N-p_1I)(D-d_2I)&=0,\\
 (D-d_0I)(D-d_1I)(N-p_2I)&=0.
\end{split}                                                   \tag{G57I.4}
\]
The right diagonal product in the first identity has only
its zeroth entry nonzero, and that entry is a unit.
Thus $N_{00}=p_0$ and $N_{10}=N_{20}=0$.  The left
diagonal product in the third identity has only its second
entry nonzero, again a unit.  Thus $N_{21}=0$ and
$N_{22}=p_2$.  The $(1,1)$ entry of the middle identity
gives $N_{11}=p_1$.  This proves the required upper
triangular form and diagonal on all inertia.  Every
cancellation here is by a unit; no integral-domain
assumption is being used.

Let now $g\in\ker\eta$, and in this basis write
$k=r(g)$ and $B=r(gsg^{-1})$.  The subgroup $I_\ell$
is normal, and the $\gamma_i$ are characters of the full
local group.  Hence $B$ is upper triangular with diagonal
$d$.  On $\ker\eta$ the matrix part of the representation
is multiplicative, so $Bk=kD$.  Its $(2,0)$ entry gives
$(d_2-d_0)k_{20}=0$, its $(2,1)$ entry gives
$(d_2-d_1)k_{21}=0$, and its $(1,0)$ entry then gives
$(d_1-d_0)k_{10}=0$.  The three differences are units,
so $k_{20}=k_{21}=k_{10}=0$.  This is upper
triangularity on the exact subgroup required by
\texttt{TriIn}.  Together with $h\equiv I$, it is
\texttt{CondT}.  At a split prime $\eta=1$ on the local
group, so it is upper triangularity on that whole group.
The matrix calculation is the integral form of
\path{CommutingUpperTriangularDistinct.lean},
\texttt{lemma2b}; the strict basis is given explicitly
by (G57I.3), rather than obtained by field descent.
\end{proof}

\paragraph{Scope and geometric meaning.}
The flag is a filtration of the actual rank-three vector
bundle over $\operatorname{Spec}A$ by free direct summands,
constructed by polynomial inertia projectors.  It is
compatible with specialization and gives the local ordinary
input on the characteristic-$\ell$ curves of
\path{RatTorsionFree.lean}, \texttt{OrdPt.ord}.
The new step is an integral local result, not a proof of
the global \texttt{FinEnd} assertion or of M2.  In
particular it supplies no finite-index subgroup on which
the characteristic polynomials of every such global curve
agree with the seed, no Type I automorphic seed, and no
numerical Selmer-length bound.
The finite local systems retain the already audited
SGA~1, Expos\'e~V, \S\S4 and~7 interpretation; the
polynomial-projector and integral flag assertions are
the explicit algebraic proof above, not an additional
consequence of that categorical dictionary.
